\documentclass[a4paper,12pt]{report}

% =============================================================
% 1. PACKAGES ET CONFIGURATION GLOBALE
% =============================================================
\usepackage[utf8]{inputenc}
\usepackage[T1]{fontenc}
\usepackage[francais,english]{babel}\selectlanguage{francais}

\usepackage{graphicx}
\usepackage{geometry}
\usepackage{xcolor}
\usepackage{fancyhdr}
\usepackage{titlesec}
\usepackage{helvet}
\usepackage{hyperref}
\usepackage{array}
\usepackage{booktabs}
\usepackage{eurosym}
\usepackage{float}
\usepackage{tikz}
\usepackage{colortbl}
\usepackage{pdflscape}
\usepackage[most]{tcolorbox}
\usepackage{enumitem}
\usepackage{longtable}
\usepackage{multirow}
\usepackage{tabularx}
\usepackage{amsmath}
\usepackage{siunitx}
\usepackage{listings}

\usetikzlibrary{shapes, arrows, positioning, shadows.blur, patterns, fit}

% Police par défaut : Sans-Serif
\renewcommand{\familydefault}{\sfdefault}
\geometry{hmargin=2.5cm, vmargin=2.5cm, headheight=45pt}

% --- PALETTE DE COULEURS (CHARTE ROBOTANK) ---
\definecolor{UTGreen}{RGB}{47, 110, 193}      % Bleu principal  #2E6EC1
\definecolor{UTSlate}{RGB}{31, 78, 121}        % Bleu foncé     #1F4E79
\definecolor{UTRed}{RGB}{192, 57, 43}          % Rouge alertes
\definecolor{UTBlue}{RGB}{22, 160, 133}        % Cyan accent
\definecolor{UTLight}{RGB}{214, 228, 240}      % Bleu très clair

% --- CONFIGURATION DES BOITES ---
\newtcolorbox{justification}[1]{
    colback=UTGreen!5!white, colframe=UTGreen, fonttitle=\bfseries,
    title=Justification Technique : #1, boxrule=0.5mm, rounded corners, drop shadow
}
\newtcolorbox{calcbox}[1]{
    colback=UTBlue!5!white, colframe=UTBlue, fonttitle=\bfseries,
    title=Calcul de Dimensionnement : #1, boxrule=0.5mm, rounded corners
}
\newtcolorbox{alertbox}[1]{
    colback=UTRed!5!white, colframe=UTRed, fonttitle=\bfseries,
    title=Analyse de Risque : #1, boxrule=0.5mm, rounded corners, drop shadow
}
\newtcolorbox{infobox}[1]{
    colback=UTSlate!5!white, colframe=UTSlate, fonttitle=\bfseries,
    title=Information : #1, boxrule=0.5mm, rounded corners
}

% --- LIENS ET METADONNÉES ---
\hypersetup{colorlinks=true, linkcolor=UTSlate, urlcolor=UTGreen, citecolor=UTGreen}

% --- COMMANDE IMAGE SÉCURISÉE ---
\newcommand{\safeincludegraphics}[2][]{%
  \IfFileExists{#2}{\includegraphics[#1]{#2}}{%
    \begin{tcolorbox}[colback=red!10, colframe=red, title=Fichier Manquant]
        \centering \textbf{Attention :} L'image \texttt{#2} est introuvable.\\
        Déposez le fichier dans le dossier du projet.
    \end{tcolorbox}%
  }%
}

% --- EN-TÊTES ---
\pagestyle{fancy}
\fancyhf{}
\lhead{\small \textbf{\color{UTSlate}SAÉ6.IOM.01 — Sujet n°10 : Robot de surveillance et de mesure}}
\rhead{\small \textbf{Cahier des Charges Technique} \\ \footnotesize Hénon Léo-Paul — IUT de Blois — vFinal}
\lfoot{\small \color{UTSlate}Léo-Paul Hénon}
\rfoot{\small \color{UTSlate}SAÉ6.IOM.01}
\cfoot{\thepage}
\renewcommand{\headrulewidth}{1pt}
\renewcommand{\headrule}{\hbox to\headwidth{\color{UTGreen}\leaders\hrule height \headrulewidth\hfill}}

% --- TITRES ---
\titleformat{\chapter}[display]
  {\normalfont\huge\bfseries\color{UTSlate}}{\chaptertitlename\ \thechapter}{20pt}{\Huge\color{UTGreen}}
\titleformat{\section}
  {\normalfont\Large\bfseries\color{UTGreen}}{\thesection}{1em}{}
\titleformat{\subsection}
  {\normalfont\large\bfseries\color{UTSlate}}{\thesubsection}{1em}{}

% --- LISTES ---
\setlist[itemize]{noitemsep, topsep=4pt}
\setlist[enumerate]{noitemsep, topsep=4pt}

% --- LISTINGS ---
\lstset{
  basicstyle=\ttfamily\small,
  breaklines=true,
  frame=single,
  rulecolor=\color{UTSlate},
  keywordstyle=\color{UTBlue}\bfseries,
  commentstyle=\color{UTSlate},
  stringstyle=\color{UTGreen},
  showstringspaces=false
}

\begin{document}

% =============================================================
% PAGE DE GARDE
% =============================================================
\begin{titlepage}
    \newgeometry{top=2cm, bottom=2cm, left=2cm, right=2cm}
    \begin{center}
        \vspace*{0.5cm}
        {\Large \textbf{\color{UTSlate}BUT Réseaux \& Télécommunications}}\\[0.2cm]
        {\small Parcours Internet et Objets Mobiles (IOM) — IUT de Blois}
        \vspace*{1.5cm}

        {\color{UTSlate}\rule{\linewidth}{1.2mm}} \\[0.5cm]
        {\Huge \bfseries \color{UTSlate} CAHIER DES CHARGES TECHNIQUE} \\[0.4cm]
        {\Large \color{UTGreen}Spécifications, Architecture \& Validation}\\[0.3cm]
        {\color{UTSlate}\rule{\linewidth}{1.2mm}}

        \vspace*{2cm}

        \begin{tcolorbox}[colback=UTLight, colframe=UTSlate, width=0.92\textwidth, arc=2mm, boxrule=0.8mm]
            \centering
            \vspace{0.3cm}
            {\Large \textbf{\color{UTSlate} ROBOTANK}}\\[0.2cm]
            {\large \textbf{\color{UTGreen}Robot de Surveillance et de Mesure}}\\[0.2cm]
            \textit{\color{UTSlate}Architecture : Raspberry Pi 4 + XIAO ESP32-C6 (Zigbee)}\\[0.1cm]
            \textit{\color{UTSlate}Caméra Fisheye IR + OLED + HC-SR04 + Steppers 28BYJ-48}
            \vspace{0.3cm}
        \end{tcolorbox}

        \vfill

        \begin{tabular}{p{0.45\textwidth} | p{0.45\textwidth}}
            \textbf{\color{UTSlate}Chef de Projet} & \textbf{\color{UTSlate}Commanditaire} \\
            \textsc{Léo-Paul Hénon} & M. Jeanneret \\
            IUT de Blois & SAÉ6.IOM.01 \\
        \end{tabular}

        \vspace*{1cm}
        {\small \color{UTSlate}Version Finale — \today}
    \end{center}
    \restoregeometry
\end{titlepage}

% =============================================================
% SUIVI DOCUMENT + SOMMAIRE
% =============================================================
\thispagestyle{empty}
\chapter*{Suivi du Document}
\begin{center}
    \begin{tabular}{|c|c|p{8cm}|c|}
    \hline
    \rowcolor{UTGreen!20} \textbf{Ver.} & \textbf{Date} & \textbf{Objet de la modification} & \textbf{Auteur} \\ \hline
    1.0 & 04/02/26 & Analyse fonctionnelle & L. Hénon \\ \hline
    2.0 & 05/02/26 & Intégration architecture hybride & L. Hénon \\ \hline
    3.0 & \today & Mise à jour RPi4 + XIAO ESP32-C6 Zigbee + steppers + charte graphique & L. Hénon \\ \hline
    \end{tabular}
\end{center}

\newpage
\tableofcontents
\newpage

% =============================================================
% CHAPITRE 0 : RAPPEL DU CDCF
% =============================================================
\chapter{Rappel du Cahier des Charges Fonctionnel (CdCF)}

\section{Description du besoin client}
Le client demande un robot mobile connecté capable de :
\begin{itemize}
    \item se déplacer en mode \textbf{autonome} ou \textbf{télécommandé} (Bluetooth),
    \item détecter un \textbf{obstacle} et adapter son comportement (arrêt / changement direction),
    \item capturer une \textbf{photo} lors d’un événement (obstacle, changement direction),
    \item mesurer la \textbf{qualité de réception Wifi} (RSSI) afin de produire une \textbf{cartographie/heatmap},
    \item fournir un \textbf{retour utilisateur local} via écran OLED et LEDs d’état.
\end{itemize}

\section{Contraintes client}
\begin{infobox}{Contraintes imposées}
\begin{itemize}
    \item utilisation \textbf{indoor} (sol plat),
    \item tension système \textbf{max 12V},
    \item autonomie minimale \textbf{20 minutes},
    \item usage prioritaire du \textbf{matériel IUT},
    \item solution \textbf{présentable} et proche d’un produit commercialisable.
\end{itemize}
\end{infobox}

\begin{infobox}{Remarque budget}
Aucune contrainte de budget n’est imposée pour ce projet (matériel fourni / stock IUT).
Le budget peut toutefois être présenté en \textit{valeur de remplacement} à titre indicatif.
\end{infobox}

% =============================================================
% CHAPITRE 1 : INTRODUCTION
% =============================================================
\chapter{Introduction et Contexte}

\section{Objectif du Projet}
Ce Cahier des Charges Technique (CdCT) encadre la réalisation du \textbf{Sujet n°10}.
Le système attendu est un robot mobile capable :
\begin{itemize}
    \item de \textbf{se déplacer} (mode autonome + mode manuel),
    \item de \textbf{surveiller} une zone (détection + capture d’image),
    \item de \textbf{mesurer} la qualité du Wifi via le RSSI,
    \item de produire une \textbf{heatmap} des mesures,
    \item d’informer l’utilisateur via \textbf{OLED + LEDs d’état}.
\end{itemize}

\section{Périmètre}
\begin{infobox}{Dans le périmètre}
\begin{itemize}
    \item Navigation indoor, évitement simple,
    \item Mesure RSSI Wifi + stockage local,
    \item Capture photo (PiCam IR 1080p),
    \item Restitution CSV + heatmap,
    \item Feedback local (OLED + LEDs).
\end{itemize}
\end{infobox}

\begin{infobox}{Hors périmètre}
\begin{itemize}
    \item SLAM/LiDAR avancé,
    \item IA “reconnaissance faciale” (non requis/interdit),
    \item temps réel garanti en zone sans Wifi.
\end{itemize}
\end{infobox}

% =============================================================
% CHAPITRE 2 : ANALYSE FONCTIONNELLE
% =============================================================
\chapter{Analyse Fonctionnelle}

\section{Bête à corne}
\begin{itemize}
    \item \textbf{À qui rend-il service ?} Technicien réseau / enseignant / démonstration
    \item \textbf{Sur quoi agit-il ?} Zone indoor + qualité Wifi
    \item \textbf{Dans quel but ?} Surveiller + cartographier la couverture Wifi
\end{itemize}

\section{Fonctions de service}
\begin{table}[H]
\centering
\rowcolors{2}{gray!10}{white}
\begin{tabular}{|c|p{11cm}|}
\hline
\rowcolor{UTSlate}\color{white}\textbf{ID} & \color{white}\textbf{Fonction principale} \\ \hline
FP1 & Se déplacer (manuel + autonome) \\ \hline
FP2 & Détecter et éviter les obstacles \\ \hline
FP3 & Mesurer le RSSI Wifi à intervalle régulier \\ \hline
FP4 & Associer une position estimée aux mesures (odométrie simple) \\ \hline
FP5 & Enregistrer les données localement (CSV/JSON) \\ \hline
FP6 & Générer une heatmap exploitable \\ \hline
FP7 & Capturer une photo lors d’un événement (obstacle/changement direction) \\ \hline
FP8 & Afficher l’état du robot (OLED + LEDs : normal/obstacle/erreur) \\ \hline
\end{tabular}
\caption{Fonctions principales (service rendu)}
\end{table}

\section{Contraintes générales}
\begin{itemize}
    \item Autonomie minimale : \textbf{20 minutes} (démo)
    \item Tension système : \textbf{max 12V}
    \item Utilisation : \textbf{indoor}
    \item Fabrication : contrainte \textbf{soudure/PCB} (shield maison)
\end{itemize}

% =============================================================
% CHAPITRE 3 : EXIGENCES
% =============================================================
\chapter{Exigences et Critères d’Acceptation}

\section{Méthode}
Chaque exigence est testable et associée à un critère de validation.

\begin{longtable}{|c|p{7.6cm}|p{3.6cm}|p{1.8cm}|}
\hline
\rowcolor{UTGreen!25}\textbf{ID} & \textbf{Exigence} & \textbf{Critère d’acceptation} & \textbf{Priorité} \\ \hline
\endhead

EX-01 & Le robot doit fonctionner au moins 20 min en autonomie. & Chrono : \(\geq\) 20 min & MUST \\ \hline
EX-02 & Le robot doit être pilotable manuellement (manette Zigbee (XIAO ESP32-C6)). & Téléopération OK 5 min & MUST \\ \hline
EX-03 & En autonome, éviter un obstacle placé à 20 cm. & 0 collision / 10 essais & MUST \\ \hline
EX-04 & Mesurer le RSSI Wifi à intervalle \(\leq 2\) s. & CSV : pas > 2s & SHOULD \\ \hline
EX-05 & Stocker localement les mesures en cas de perte Wifi. & logs présents offline & MUST \\ \hline
EX-06 & Produire une heatmap exportable (PNG). & fichier PNG lisible & SHOULD \\ \hline
EX-07 & Capturer une photo horodatée lors d’un événement. & image horodatée & SHOULD \\ \hline
EX-08 & Disposer d’un arrêt d’urgence logiciel. & STOP immédiat & MUST \\ \hline
EX-09 & La tension 5V du Raspberry Pi doit rester stable (4.9–5.2V). & mesure stable & MUST \\ \hline
EX-10 & Afficher l’état sur OLED (mode + RSSI + batterie si dispo). & info visible & SHOULD \\ \hline
EX-11 & LED verte = fonctionnement normal. & allumée en RUN & SHOULD \\ \hline
EX-12 & LED rouge = obstacle détecté. & allumée si obstacle & MUST \\ \hline
EX-13 & LED orange = erreur / sous-tension / défaut capteur. & allumée si défaut & SHOULD \\ \hline

\caption{Exigences testables et critères d'acceptation}
\end{longtable}

\begin{justification}{Priorisation}
\textbf{MUST} = nécessaire pour valider la démo.  
\textbf{SHOULD} = améliore le rendu et l’expérience utilisateur.
\end{justification}

% =============================================================
% CHAPITRE 4 : SOLUTIONS TECHNIQUES & MATERIEL
% =============================================================
\chapter{Solutions Techniques \& Matériel}

\section{Justification de l'Architecture Hybride}
\begin{justification}{Hybride RPi + Kit Arduino}
Le kit Arduino Engineering fournit une base mécanique (moteurs + châssis).  
Le Raspberry Pi apporte la puissance de calcul, le Wifi, la caméra IR, l’affichage OLED et l’intégration logicielle Python.
\end{justification}
\newpage
\section{Liste exhaustive des composants (BOM)}
\begin{infobox}{Principe}
La liste ci-dessous constitue la liste exhaustive des composants nécessaires au fonctionnement du robot (matériel principal + interface + câblage + sécurité).
\end{infobox}

\begin{longtable}{|p{6cm}|c|p{6cm}|}
\hline
\rowcolor{UTSlate}\color{white}\textbf{Composant} & \color{white}\textbf{Qté} & \color{white}\textbf{Rôle / Remarques} \\ \hline
\endhead

\textbf{Raspberry Pi 4 Model B 4Go} & 1 & unité centrale, Wifi, traitements, logs \\ \hline
\textbf{Carte microSD 16/32Go} & 1 & OS + stockage logs \\ \hline
\textbf{Raspberry Pi Camera IR 1080p} & 1 & photo/vision nocturne IR \\ \hline
\textbf{Capteur ultrason HC-SR04 ou SR04M} & 1 & mesure distance obstacle \\ \hline
\textbf{Écran OLED I2C 0.96'' (SSD1306)} & 1 & affichage mode, RSSI, état \\ \hline
\textbf{LED Verte 5mm} & 1 & état RUN \\ \hline
\textbf{LED Rouge 5mm} & 1 & obstacle / stop sécurité \\ \hline
\textbf{LED Orange 5mm} & 1 & défaut capteur / erreur \\ \hline
\textbf{Résistances 220$\Omega$} & 3 & limitation courant LEDs \\ \hline
\textbf{Adaptation niveau ECHO (diviseur)} & 2 & ex: 1k/2k ou 10k/20k (5V$\rightarrow$3.3V) \\ \hline
\textbf{Driver ULN2003 (x2)} & 1 & commande moteurs DC \\ \hline
\textbf{Châssis 3D imprimé + Moteurs stepper 28BYJ-48 + Roues} & 1 & base mécanique \\ \hline
\textbf{Batterie LiPo 2S 7.4V} & 1 & alimentation \\ \hline
\textbf{Chargeur LiPo adapté} & 1 & sécurité charge \\ \hline
\textbf{Convertisseur DC-DC 12V$\rightarrow$5V (3A min)} & 1 & alimentation stable du RPi \\ \hline
\textbf{Condensateurs (1000µF + 100nF)} & 2+ & découplage anti-reboot \\ \hline
\textbf{Interrupteur général} & 1 & coupure alimentation \\ \hline
\textbf{Fusible / protection (option)} & 1 & sécurité \\ \hline
\textbf{Breadboard / plaque d’essai} & 1 & prototypage \\ \hline
\textbf{PCB vierge (shield maison)} & 1 & version soudée \\ \hline
\textbf{Câbles Dupont M/F} & 1 lot & connexions prototypes \\ \hline
\textbf{Connecteurs JST / borniers} & 1 lot & connexions batterie/moteurs \\ \hline
\textbf{Colliers + gaine thermo} & 1 lot & finition, sécurité \\ \hline

\caption{BOM (Bill Of Materials) – liste exhaustive des composants}
\end{longtable}

\section{Capteurs \& Interface utilisateur}
\begin{itemize}
    \item \textbf{Ultrason} : obstacle (distance), déclenche LED rouge + photo si événement.
    \item \textbf{PiCam IR} : capture 1080p, utile même faible luminosité.
    \item \textbf{OLED} : affiche mode (manuel/autonome), RSSI, état.
    \item \textbf{LEDs} : feedback instantané (RUN / obstacle / défaut).
\end{itemize}

% =============================================================
% CHAPITRE 5 : ARCHITECTURE SYSTEME
% =============================================================
\chapter{Architecture Système}

\section{Schéma synoptique}
\begin{figure}[H]
\centering
\begin{tikzpicture}[
    scale=0.7, transform shape,
    node distance=3cm,
    block/.style={rectangle, draw=UTSlate, fill=white, text width=3.4cm, text centered, rounded corners, minimum height=1.5cm, drop shadow},
    rpi/.style={rectangle, draw=UTGreen, fill=UTGreen!10, text width=4.8cm, text centered, minimum height=3cm, line width=1.5pt, drop shadow},
    line_data/.style={draw, -latex, thick, color=UTBlue},
    line_power/.style={draw, -latex, thick, color=UTRed, dashed}
]
% --- NOEUDS ---
\node [rpi] (cpu) {\textbf{Raspberry Pi 4}\\(Master)};
\node [block, left=6cm of cpu] (sensors) {\textbf{Capteurs}\\HC-SR04 + XIAO\#2};
\node [block, right=6cm of cpu] (actuators) {\textbf{Steppers 28BYJ-48}\\Driver ULN2003};
\node [block, above=2.6cm of cpu] (cam) {\textbf{PiCam IR 1080p}};
\node [block, below=2.6cm of cpu] (ui) {\textbf{OLED I2C}\\+ LEDs};
\node [block, below=2.8cm of ui] (power) {\textbf{Batterie LiPo 2S}\\7.4V + DC-DC 5V};

% --- LIGNES DONNEES ---
\draw[line_data] (sensors) -- node[above, yshift=2pt] {\small GPIO} (cpu);
\draw[line_data] (cam) -- node[right, xshift=3pt] {\small CSI} (cpu);
\draw[line_data] (cpu) -- node[above, yshift=2pt] {\small PWM/GPIO} (actuators);
\draw[line_data] (cpu) -- node[right, xshift=3pt] {\small I2C/GPIO} (ui);

% --- LIGNES PUISSANCE ---
% 11.1V vers driver moteurs
\draw[line_power] (power.east) -| node[pos=0.75, below, xshift=10pt] {\small 7.4V} (actuators.south);

% 5V stable vers Raspberry Pi
\draw[line_power] (power.north) -- node[right, xshift=3pt] {\small 5V stable} (cpu.south);

% (optionnel mais logique) 5V vers OLED+LEDs (si tu veux le montrer)
\draw[line_power] (cpu.south) -- (ui.north);

% --- LEGENDE ---
\node[draw, align=left] at (0, -7.2) {\textcolor{UTBlue}{\textbf{---}} Données \quad \textcolor{UTRed}{\textbf{- - -}} Puissance};

\end{tikzpicture}
\caption{Architecture matérielle RoboTank (RPi4 + Zigbee + Steppers)}
\end{figure}

% =============================================================
% CHAPITRE 6 : ARCHITECTURE LOGICIELLE
% =============================================================
\chapter{Architecture Logicielle}

\section{Modules (organisation)}
\begin{itemize}
    \item \textbf{server.py} : serveur Flask principal, API REST, threads dédiés
    \item \textbf{xiao\_reader()} : lecture /dev/ttyACM0, commandes Zigbee F/B/L/R/S
    \item \textbf{explore\_loop()} : dead reckoning, marquage grille, mesure RSSI
    \item \textbf{\_cam\_loop()} : capture Picamera2 MJPEG 30fps
    \item \textbf{avoidance\_loop()} : évitement obstacles HC-SR04 + rotation 90°/180°
    \item \textbf{serial\_reader()} : lecture HC-SR04 via /dev/ttyUSB0
    \item \textbf{index.html} : interface web Flask (joystick, carte, caméra, vitesse)
    \item \textbf{/api/map} : grille JSON 100×100 + position robot (dead reckoning)
\end{itemize}

\section{Flux de données}
\begin{itemize}
    \item Joystick $\rightarrow$ XIAO\#1 (Zigbee ED) $\rightarrow$ XIAO\#2 (Coordinator) $\rightarrow$ /dev/ttyACM0 $\rightarrow$ drive()
    \item HC-SR04 $\rightarrow$ /dev/ttyUSB0 $\rightarrow$ avoidance\_loop() $\rightarrow$ rotation automatique
    \item explore\_loop() $\rightarrow$ dead reckoning (x,y,$\theta$) $\rightarrow$ grille JSON $\rightarrow$ canvas web
    \item Picamera2 $\rightarrow$ MJPEG $\rightarrow$ /video\_feed $\rightarrow$ interface web
\end{itemize}

% =============================================================
% CHAPITRE 7 : DIMENSIONNEMENT ELECTRIQUE
% =============================================================
\chapter{Dimensionnement \& Énergie}

\section{Hypothèses de consommation (mise à jour)}
\begin{table}[H]
\centering
\begin{tabular}{|p{6cm}|c|c|}
\hline
\rowcolor{UTGreen!25}\textbf{Élément} & \textbf{Puissance} & \textbf{Remarques} \\ \hline
Raspberry Pi 4 Model B 4Go & 3.5W & charge moyenne \\ \hline
PiCam IR 1080p & 1.0W & capture ponctuelle \\ \hline
OLED I2C & 0.3W & affichage continu \\ \hline
LEDs (3x) & 0.2W & selon état \\ \hline
Steppers 28BYJ-48 (x2) + ULN2003 & 3.0W & moyenne, hors pointe \\ \hline
Ultrason & 0.2W & négligeable \\ \hline
\textbf{Total} & \textbf{7.7W} & estimation moyenne \\ \hline
\end{tabular}
\caption{Bilan de puissance estimé (avec OLED + LEDs + PiCam)}
\end{table}

\section{Courant moyen}
\[
I_{moy}=\frac{P}{U}=\frac{7.7}{11.1}\approx 0.69A
\]

\section{Autonomie}
\begin{calcbox}{Autonomie théorique}
Énergie batterie : \(11.1V \times 0.8Ah = 8.88Wh\).\\
Rendement régulation estimé : \(0.9\).\\[2mm]
\[
T=\frac{8.88\times 0.9}{7.7}\approx 1.04h \approx 1h02
\]
\textbf{Conclusion :} autonomie largement suffisante pour 20 minutes de démonstration.
\end{calcbox}

\section{Sécurité batterie LiPo}
\begin{alertbox}{Risque LiPo}
\textbf{Risque :} surcharge, décharge profonde, gonflement.\\
\textbf{Mesures :} chargeur adapté, stockage à 3.8V/cell, inspection visuelle, interrupteur général.
\end{alertbox}

% =============================================================
% CHAPITRE 8 : RISQUES
% =============================================================
\chapter{Analyse des Risques}

\begin{alertbox}{Perte Wifi}
\textbf{Risque :} RSSI faible / trous de couverture.\\
\textbf{Mesure :} logs offline + export différé.
\end{alertbox}

\begin{alertbox}{Collision / blocage}
\textbf{Risque :} obstacle non détecté / faux négatif.\\
\textbf{Mesures :} ultrason prioritaire + stop + LED rouge + réduction vitesse.
\end{alertbox}

\begin{alertbox}{Sous-tension 5V}
\textbf{Risque :} reboot Raspberry Pi lors des pointes moteur.\\
\textbf{Mesures :} DC-DC 5V 3A min, masses communes, câblage court, condensateurs.
\end{alertbox}

% =============================================================
% CHAPITRE 9 : CYBERSECURITE
% =============================================================
\chapter{Sécurité et Cybersécurité}

\section{Menaces}
\begin{itemize}
    \item accès non autorisé (SSH)
    \item interception Wifi
    \item extraction données si SD accessible
\end{itemize}

\section{Mesures}
\begin{itemize}
    \item identifiants changés, SSH par clé
    \item pare-feu UFW (ports minimaux)
    \item données non sensibles (RSSI/position)
\end{itemize}

% =============================================================
% CHAPITRE 10 : VALIDATION
% =============================================================
\chapter{Plan de Tests et Validation}

\section{Tests unitaires}
\begin{itemize}
    \item Test moteurs : marche avant/arrière, PWM
    \item Test ultrason : distance OK à 10/20/50cm
    \item Test Wifi : RSSI loggé
    \item Test caméra : photo horodatée (jour + pénombre)
    \item Test OLED : affichage mode + RSSI
    \item Test LEDs : vert RUN, rouge obstacle, orange défaut
\end{itemize}

\section{Recette système}
\begin{table}[H]
\centering
\rowcolors{2}{gray!10}{white}
\begin{tabular}{|c|p{7cm}|p{4cm}|}
\hline
\rowcolor{UTSlate}\color{white}\textbf{ID} & \color{white}\textbf{Test} & \color{white}\textbf{Résultat attendu} \\ \hline
T-01 & Autonomie 20 min & OK \(\geq\) 20 min \\ \hline
T-02 & Évitement obstacle 20 cm (10 essais) & 0 collision \\ \hline
T-03 & Photo sur événement obstacle & image horodatée \\ \hline
T-04 & OLED affiche mode + RSSI & lisible \\ \hline
T-05 & LED rouge obstacle & allumée \\ \hline
T-06 & Heatmap générée après parcours & PNG valide \\ \hline
T-07 & Mode offline Wifi coupé & logs présents \\ \hline
\end{tabular}
\caption{Table de recette (mise à jour)}
\end{table}

% =============================================================
% CHAPITRE 11 : PLANIFICATION
% =============================================================
\chapter{Planification}

\section{Diagramme de Gantt}
\begin{landscape}
    \thispagestyle{empty}
    \begin{figure}[H]
        \centering
        \vspace*{1cm}
        \begin{tcolorbox}[colback=UTLight,colframe=UTSlate,title=Diagramme de Gantt]\centering\textit{Insérer ici le diagramme de Gantt (fichier gantt.png)}\end{tcolorbox}
        \caption{Planning Gantt du Projet (Livrable 1)}
    \end{figure}
\end{landscape}

% =============================================================
% CHAPITRE 12 : NOTICE UTILISATION (OBLIGATOIRE)
% =============================================================
\chapter{Notice d’Utilisation}

\section{Mise sous tension}
\begin{enumerate}
    \item Vérifier la batterie et les connexions (moteurs + DC-DC + RPi).
    \item Actionner l’interrupteur général.
    \item Attendre le démarrage du Raspberry Pi (\~30s).
\end{enumerate}

\section{Signification OLED + LEDs}
\begin{itemize}
    \item \textbf{LED verte} : fonctionnement normal (RUN)
    \item \textbf{LED rouge} : obstacle détecté / arrêt sécurité
    \item \textbf{LED orange} : défaut système (capteur/alim)
\end{itemize}

\section{Modes}
\begin{itemize}
    \item \textbf{Mode manuel} : pilotage via manette Zigbee physique
    \item \textbf{Mode autonome} : déplacement + évitement + mesures RSSI
\end{itemize}

\section{Récupération des données}
Les logs et photos sont stockés sur la carte SD :
\begin{itemize}
    \item \texttt{/home/pi/logs/} : CSV/JSON
    \item \texttt{/home/pi/photos/} : images horodatées
\end{itemize}

% =============================================================
% CHAPITRE 13 : NOTICE MISE À JOUR / MAINTENANCE (OBLIGATOIRE)
% =============================================================
\chapter{Notice de Maintenance et Mise à Jour}

\section{Maintenance préventive}
\begin{itemize}
    \item inspection batterie LiPo (gonflement)
    \item vérification serrage mécanique
    \item contrôle câbles moteurs et connecteurs
    \item nettoyage optique caméra
\end{itemize}

\section{Mise à jour logicielle (exemple)}
\begin{lstlisting}[language=bash]
cd /home/pi/robot
git pull
sudo reboot
\end{lstlisting}

% =============================================================
% CONCLUSION
% =============================================================
\chapter{Conclusion}

Ce cahier des charges technique formalise les exigences, l’architecture et la validation du robot.
L’ajout de la PiCam IR, de l’OLED et des LEDs améliore fortement l’expérience utilisateur, la surveillance et la robustesse en démonstration.
Le système répond aux fonctionnalités demandées : déplacement, évitement, mesure RSSI et cartographie.

% =============================================================
% ANNEXES
% =============================================================
\appendix
\chapter{Annexes}

\section{Fichiers attendus dans le dossier}
\begin{itemize}
    \item \texttt{logo\_iut\_blois.png}
    \item \texttt{gantt.png}
\end{itemize}

\end{document}